\documentclass[a4paper,12pt]{article}
\usepackage[T2A]{fontenc}
\usepackage[utf8]{inputenc}
\usepackage[russian]{babel}
\usepackage{amsmath,amssymb,mathtools}
\usepackage{PTSans}
\renewcommand{\familydefault}{\sfdefault}
\usepackage{icomma}
\usepackage[a4paper,margin=18mm,headheight=25pt]{geometry}
\usepackage{microtype}
\usepackage{tikz}
\usetikzlibrary{calc,arrows.meta,positioning,shapes.geometric}
\usepackage{tabularx,booktabs,longtable}
\usepackage{enumitem}
\usepackage{hyperref}
\usepackage{parskip}
\usepackage{fancyhdr}
\usepackage{setspace}
\usepackage{xcolor}

\definecolor{accent}{HTML}{2D6F6D}
\definecolor{darkaccent}{HTML}{1D4D4B}
\definecolor{textgray}{HTML}{2F3437}
\definecolor{softgray}{HTML}{6B7478}
\definecolor{linegray}{HTML}{D9DEDF}

\hypersetup{colorlinks=true,linkcolor=darkaccent,urlcolor=darkaccent}
\setlength{\parindent}{0pt}
\setlength{\parskip}{5pt}
\onehalfspacing
\setlist[itemize]{leftmargin=1.4em,itemsep=2pt,topsep=3pt}
\setlist[enumerate]{leftmargin=1.8em,itemsep=5pt,topsep=4pt}

\pagestyle{fancy}
\fancyhf{}
\fancyhead[L]{\small\color{softgray}Екатерина Скелли}
\fancyhead[R]{\small\color{softgray}03.10.26}
\fancyfoot[C]{\small\color{softgray}\thepage}
\renewcommand{\headrulewidth}{0.3pt}
\renewcommand{\headrule}{\hbox to\headwidth{\color{linegray}\leaders\hrule height \headrulewidth\hfill}}

\newcommand{\sectiontitle}[1]{%
  \vspace{2mm}%
  {\Large\bfseries\color{darkaccent}#1}\par\vspace{1mm}%
  {\color{linegray}\hrule}\vspace{2mm}%
}
\newcommand{\key}[1]{\textbf{\color{darkaccent}#1}}
\newcommand{\warn}[1]{\textbf{\color{accent}Важно:} #1}

\tikzset{
  flowstart/.style={ellipse,draw=accent,very thick,align=center,minimum width=42mm,minimum height=10mm,fill=accent!6},
  flowaction/.style={rounded corners=3pt,rectangle,draw=accent,very thick,align=center,text width=84mm,minimum height=12mm,fill=accent!4},
  flowdecision/.style={diamond,aspect=2.2,draw=accent,very thick,align=center,text width=48mm,inner sep=1.5mm,fill=accent!5},
  flowarrow/.style={-{Stealth[length=3mm]},very thick,draw=darkaccent}
}

\begin{document}

% ---------------- Title page ----------------
\thispagestyle{empty}
\begin{center}
\vspace*{25mm}
{\small\color{softgray}МАТЕМАТИКА}\par
\vspace{8mm}
{\Huge\bfseries\color{darkaccent}Чек-лист}\par
\vspace{5mm}
{\LARGE\bfseries Числовые промежутки, свойства неравенств\\[2mm]
и движение протяжённых тел}\par
\vspace{12mm}
{\large 3 октября 2026 г.}\par
\vfill
{\large\bfseries Лёвин Артём Александрович}\par
\vspace{2mm}
{\large эксклюзивно для Екатерины Скелли}\par
\vspace{20mm}
\begin{tikzpicture}[x=1cm,y=1cm]
  \draw[accent!75, line width=1.15pt]
    (0,0) .. controls (2.2,0.8) and (4.4,-0.7) .. (6.7,0.05)
          .. controls (8.9,0.75) and (11.2,-0.35) .. (14,0.35);
  \draw[accent!25, line width=0.7pt]
    (1,-0.35) .. controls (3.0,0.15) and (5.0,-0.55) .. (7.2,-0.1)
          .. controls (9.1,0.25) and (11.5,-0.65) .. (13.2,-0.05);
\end{tikzpicture}
\end{center}
\clearpage

% ---------------- Main checklist: intervals ----------------
\sectiontitle{1. Числовые промежутки: как читать и записывать}

\key{Главная идея.} Промежуток --- это множество чисел на числовой прямой. Вид скобки показывает, входит ли граничная точка в множество.

\begin{center}
\renewcommand{\arraystretch}{1.35}
\begin{tabularx}{0.92\linewidth}{@{}lXc@{}}
\toprule
Запись & Что означает & Точка на прямой \\
\midrule
$x<a$ или $x>a$ & граница не входит & выколотая \\
$x\le a$ или $x\ge a$ & граница входит & закрашенная \\
$(a;b)$ & обе границы не входят & две выколотые \\
$[a;b]$ & обе границы входят & две закрашенные \\
$[a;b)$, $(a;b]$ & входит только одна граница & смешанный случай \\
\bottomrule
\end{tabularx}
\end{center}

\warn{Круглая скобка соответствует строгому неравенству; квадратная --- нестрогому.}

\textbf{Пример.} Неравенство $-5\le x<-3$ записывается как
\[
 x\in[-5;-3).
\]

\begin{center}
\begin{tikzpicture}[x=1.25cm,y=1cm]
  \draw[-{Stealth[length=2.5mm]},line width=0.8pt] (-6.2,0)--(-1.6,0);
  \foreach \x in {-6,-5,-4,-3,-2}{
    \draw (\x,0.09)--(\x,-0.09) node[below=2mm]{\small \x};
  }
  \draw[accent,line width=3pt] (-5,0)--(-3,0);
  \fill[accent] (-5,0) circle (2.4pt);
  \draw[accent,line width=1.2pt,fill=white] (-3,0) circle (2.7pt);
\end{tikzpicture}
\end{center}

\subsection*{Пересечение и объединение}
\begin{itemize}
  \item \key{Пересечение $A\cap B$} --- числа, которые принадлежат обоим промежуткам одновременно.
  \item \key{Объединение $A\cup B$} --- числа, которые принадлежат хотя бы одному из промежутков.
\end{itemize}

\textbf{Пример.} Пусть $A=(2;4)$ и $B=[0;3]$. Тогда
\[
A\cap B=(2;3],\qquad A\cup B=[0;4).
\]

\subsection*{Практическая задача: температура и промежуток времени}
В 7:00 температура равна $16^\circ\!\mathrm{C}$ и растёт на $2^\circ\!\mathrm{C}$ в час. В 14:00 она достигает $30^\circ\!\mathrm{C}$, после чего падает на $3^\circ\!\mathrm{C}$ в час.

Температура впервые становится равной $20^\circ\!\mathrm{C}$ в 9:00. После 14:00 ей нужно потерять $10^\circ\!\mathrm{C}$, то есть
\[
\frac{10}{3}\text{ ч}=3\text{ ч }20\text{ мин}.
\]
Значит, температура не ниже $20^\circ\!\mathrm{C}$ на промежутке
\[
\boxed{[9{:}00;17{:}20]}.
\]

\subsection*{Сравнение дробей с десятичным промежутком}
Если нужно выяснить, какая дробь лежит между $0,5$ и $0,6$, можно:
\begin{itemize}
  \item перевести дробь в десятичную;
  \item привести границы к тому же знаменателю;
  \item использовать грубую оценку и исключение невозможных вариантов.
\end{itemize}
Например,
\[
\frac{5}{9}=0,(5),\qquad 0,5<\frac{5}{9}<0,6.
\]

\key{Полезное наблюдение.} Если $0<m<1$, то
\[
m^2<m<m+1.
\]
При возведении положительного числа меньше единицы в квадрат оно уменьшается.

% ---------------- Inequalities ----------------
\sectiontitle{2. Свойства неравенств}

Пусть известно, что $a<b$. Тогда:

\begin{center}
\renewcommand{\arraystretch}{1.45}
\begin{tabularx}{0.94\linewidth}{@{}XcX@{}}
\toprule
Действие & Знак & Результат \\
\midrule
Прибавить или вычесть одно и то же число & сохраняется & $a+c<b+c$ \\
Умножить или разделить на положительное число & сохраняется & $ac<bc$, если $c>0$ \\
Умножить или разделить на отрицательное число & меняется & $ac>bc$, если $c<0$ \\
\bottomrule
\end{tabularx}
\end{center}

\warn{Смена знака происходит именно при умножении или делении обеих частей на отрицательное число.}

\subsection*{Как сравнивать преобразованные выражения}
Если $m<n$, то:
\[
m+9<n+9,
\]
\[
-20m>-20n,
\]
а для сравнения $4-m$ и $4-n$ удобно сделать два шага:
\[
m<n\Rightarrow -m>-n\Rightarrow 4-m>4-n.
\]

\subsection*{Как оценивать выражение}
\textbf{Пример.} Известно $b>3$. Оценим выражение $\dfrac{b}{-2}+10$.

Делим неравенство на $-2$ и меняем знак:
\[
\frac{b}{-2}<-\frac32.
\]
Прибавляем $10$ к обеим частям:
\[
\frac{b}{-2}+10<8,5.
\]

\subsection*{Сравнение через заданное равенство}
Пусть
\[
a+1=2b,\qquad b>1.
\]
Тогда $a=2b-1$, поэтому
\[
a-b=(2b-1)-b=b-1>0.
\]
Следовательно,
\[
\boxed{a>b}.
\]
Это строгий способ доказать сравнение: свести вопрос к знаку разности $a-b$.

\subsection*{Почленное сложение}
Если одновременно $x>y$ и $p>q$, то можно сложить левые и правые части:
\[
x+p>y+q.
\]
Используйте это только тогда, когда направления исходных неравенств согласованы.

\clearpage

% ---------------- Flowchart ----------------
\thispagestyle{fancy}
\begin{center}
{\LARGE\bfseries\color{darkaccent}Алгоритм: оценка выражения по неравенству}\par
\vspace{12mm}
\begin{tikzpicture}[node distance=10mm,
  flowaction/.append style={text width=68mm},
  flowdecision/.append style={text width=42mm}]
  \node[flowstart] (s) {Дано базовое неравенство};
  \node[flowaction,below=of s] (target) {Определите, какое выражение нужно получить};
  \node[flowdecision,below=13mm of target] (mult) {Нужно сначала умножить или разделить?};
  \node[flowdecision,below=14mm of mult] (sign) {Множитель или делитель отрицательный?};
  \node[flowaction,below left=13mm and 10mm of sign,text width=53mm] (flip) {Поменяйте направление знака};
  \node[flowaction,below right=13mm and 10mm of sign,text width=53mm] (same) {Сохраните направление знака};
  \node[flowaction,below=30mm of sign] (add) {Прибавьте или вычтите нужное число из обеих частей; знак сохраняется};
  \node[flowaction,below=of add] (check) {Сверьте левую часть с требуемым выражением и упростите правую};
  \node[flowstart,below=of check] (end) {Запишите итоговую оценку};

  \draw[flowarrow] (s)--(target);
  \draw[flowarrow] (target)--(mult);
  \draw[flowarrow] (mult)-- node[right,fill=white,inner sep=1pt]{\small да} (sign);
  \draw[flowarrow] (mult.east) -| node[pos=0.25,above,fill=white,inner sep=1pt]{\small нет} (add.east);
  \draw[flowarrow] (sign)-- node[above left,fill=white,inner sep=1pt]{\small да} (flip);
  \draw[flowarrow] (sign)-- node[above right,fill=white,inner sep=1pt]{\small нет} (same);
  \draw[flowarrow] (flip.south) |- (add.west);
  \draw[flowarrow] (same.south) |- (add.east);
  \draw[flowarrow] (add)--(check);
  \draw[flowarrow] (check)--(end);
\end{tikzpicture}
\end{center}
\clearpage

% ---------------- Motion ----------------
\sectiontitle{3. Движение протяжённых тел}

\key{Главная идея.} У поезда, баржи или другого длинного объекта нельзя всегда считать только путь одной точки. Для полного прохождения мимо другого объекта нужно учесть протяжённости и относительное движение.

\subsection*{1. Скорость сближения}
\begin{itemize}
  \item навстречу друг другу: $v_{\text{отн}}=v_1+v_2$;
  \item в одном направлении: $v_{\text{отн}}=|v_1-v_2|$;
  \item неподвижный объект: его скорость равна $0$.
\end{itemize}

\subsection*{2. Относительный путь}
Для полного прохождения одного протяжённого объекта мимо другого:
\[
S_{\text{отн}}=L_1+L_2.
\]
Если до начала обгона есть отставание $d_1$, а после обгона нужно получить опережение $d_2$, то
\[
S_{\text{отн}}=L_1+L_2+d_1+d_2.
\]

\subsection*{3. Базовая формула}
\[
\boxed{t=\frac{S_{\text{отн}}}{v_{\text{отн}}}}
\]
или, если неизвестен относительный путь,
\[
S_{\text{отн}}=v_{\text{отн}}t.
\]

\warn{Все величины должны быть согласованы по единицам. Если скорость дана в км/ч, длины удобно переводить в километры, а время --- в часы.}

\begin{center}
\begin{tikzpicture}[x=0.92cm,y=1cm]
  \draw[->,line width=0.8pt] (0,0)--(13,0);
  \node[above] at (11.1,0.05){\small направление движения};
  \draw[fill=accent!10,draw=accent,very thick] (0.8,0.45) rectangle (4.6,1.25);
  \node at (2.7,0.85){\small объект 1, длина $L_1$};
  \draw[fill=accent!5,draw=accent,very thick] (7.4,0.45) rectangle (10.4,1.25);
  \node at (8.9,0.85){\small объект 2, длина $L_2$};
  \draw[<->,darkaccent,line width=1pt] (4.6,1.55)--(7.4,1.55) node[midway,above]{\small отставание $d_1$};
\end{tikzpicture}
\end{center}

\subsection*{Пример 1. Поезд и лесополоса}
Поезд движется со скоростью $90$ км/ч и за $1$ минуту полностью проходит вдоль лесополосы длиной $800$ м. Найдём длину поезда $x$.

\[
1\text{ мин}=\frac1{60}\text{ ч},\qquad 800\text{ м}=0,8\text{ км}.
\]
\[
\frac1{60}=\frac{x+0,8}{90}.
\]
Отсюда
\[
x+0,8=1,5,\qquad x=0,7\text{ км}=700\text{ м}.
\]

\subsection*{Пример 2. Два движущихся объекта}
Если поезд со скоростью $65$ км/ч обгоняет пешехода, идущего в том же направлении со скоростью $5$ км/ч, то
\[
v_{\text{отн}}=65-5=60\text{ км/ч}.
\]
За $30$ секунд относительный путь равен
\[
60\cdot\frac{30}{3600}=0,5\text{ км}=500\text{ м}.
\]
Если длиной пешехода пренебрегаем, это и есть длина поезда.

\subsection*{Пример 3. Догнал, прошёл мимо, перегнал}
Длины двух барж равны $60$ м и $40$ м. Первая баржа вначале отстаёт на $200$ м, а после полного обгона её нос находится на $300$ м впереди второй. На весь манёвр ушло $18$ минут.

Полный относительный путь:
\[
60+40+200+300=600\text{ м}=0,6\text{ км}.
\]
Время:
\[
18\text{ мин}=0,3\text{ ч}.
\]
Поэтому разность скоростей равна
\[
\boxed{v_1-v_2=\frac{0,6}{0,3}=2\text{ км/ч}}.
\]

\subsection*{Мини-памятка}
\begin{enumerate}
  \item Определите, какие объекты имеют ненулевую длину.
  \item Определите направление движения и найдите $v_{\text{отн}}$.
  \item Составьте полный относительный путь: длины объектов, отставание, опережение.
  \item Приведите единицы к одной системе.
  \item Используйте $t=S_{\text{отн}}/v_{\text{отн}}$ и проверьте смысл ответа.
\end{enumerate}

\clearpage

% ---------------- Training ----------------
\sectiontitle{Тренировочный блок --- Путь А}

\textbf{1.} Запишите промежутком множество решений неравенства
\[
-4\le x<2.
\]

\textbf{2.} Даны промежутки $A=(-3;5]$ и $B=[1;7)$. Найдите $A\cap B$ и $A\cup B$.

\textbf{3.} Какие из чисел принадлежат промежутку $(0,4;0,6)$?
\[
\frac25,\qquad \frac12,\qquad \frac35,\qquad \frac59.
\]

\textbf{4.} В 8:00 температура была $15^\circ\!\mathrm{C}$. До 12:00 она равномерно росла на $3^\circ\!\mathrm{C}$ в час, затем стала равномерно снижаться на $2^\circ\!\mathrm{C}$ в час. Укажите промежуток времени, когда температура была не ниже $21^\circ\!\mathrm{C}$.

\textbf{5.} Известно, что $p<q$. Поставьте знак $<$ или $>$:
\[
p+6\ \square\ q+6,\qquad -3p\ \square\ -3q,\qquad 5-p\ \square\ 5-q.
\]

\textbf{6.} Известно, что $a>12$. Оцените выражение
\[
-0,5a+4.
\]

\textbf{7.} Известно, что
\[
a+3=2b,\qquad b>3.
\]
Сравните $a$ и $b$.

\textbf{8.} Поезд движется со скоростью $72$ км/ч и проходит мимо неподвижного столба за $25$ секунд. Найдите длину поезда в метрах.

\textbf{9.} Поезд движется со скоростью $90$ км/ч и полностью проходит мост длиной $600$ м за $48$ секунд. Найдите длину поезда.

\textbf{10.} Длины двух барж равны $60$ м и $40$ м. Первая баржа движется быстрее второй в том же направлении. Вначале она отстаёт на $200$ м, а после полного обгона оказывается на $300$ м впереди. Весь манёвр занял $18$ минут. На сколько километров в час скорость первой баржи больше скорости второй?

\clearpage

% ---------------- Answers ----------------
\sectiontitle{Ответы}

\begin{center}
\renewcommand{\arraystretch}{1.45}
\begin{tabularx}{0.96\linewidth}{@{}cX@{}}
\toprule
№ & Ответ \\
\midrule
1 & $[-4;2)$ \\
2 & $A\cap B=[1;5]$, \quad $A\cup B=(-3;7)$ \\
3 & $\dfrac12$ и $\dfrac59$ \\
4 & $[10{:}00;15{:}00]$ \\
5 & $<$, $>$, $>$ \\
6 & $-0,5a+4<-2$ \\
7 & $a>b$ \\
8 & $500$ м \\
9 & $600$ м \\
10 & $2$ км/ч \\
\bottomrule
\end{tabularx}
\end{center}

\vfill
{\small\color{softgray}Лёвин Артём Александрович эксклюзивно для Екатерины Скелли}

\end{document}
